%Usual structure:
% Background and Motivation;
% Contextualization;
% Problem or Focus;
% Research Objective/Question
% Methodology
% Tasks and Schedule
% Main Contributions and Expected Results
% Dissertation Structure


\chapter{Introduction} \label{chap:intro}

TODO: Reescrever introdução com as devidas alterações


Almost everything in life depends on negotiation. From the moment we are born, we negotiate for attention and care. Developed in part through the unconscious imitation of those around us \cite{Graham2014}, negotiation requires the ability to infer others’ intentions and to act under uncertainty.

Large language models (LLMs) have demonstrated impressive capabilities across a number of tasks, from code generation to creative writing and question answering. As these models are increasingly deployed in agentic scenarios—where they must interact with other agents or humans to accomplish goals—negotiation appears as a particularly interesting benchmark.

\textcite{bianchi_how_2024} introduced NegotiationArena, an open-source framework for evaluating the negotiation abilities of LLM agents across scenarios such as the Ultimatum Game, resource trading, and price negotiation. Although highly valuable, the original evaluation was limited to closed-source models (GPT-4, GPT-3.5, Claude 2.1, and Claude 2), some of which are no longer available. This limitation raises concerns regarding reproducibility and accessibility.

Techniques such as Chain-of-Thought (CoT) allow models to think step-by-step, dividing complex problems into manageable sub-problems to reduce errors and improve interpretability. Similarly, Self-Correction mechanisms (or Self-Refine) enable models to critique and improve their own outputs through an iterative process of drafting, feedback, and refinement. While these emergent abilities have demonstrated success in domains like mathematics and coding , their specific impact on negotiation has not yet been fully studied. 

Beyond individual reasoning, multi-agent collaboration offers another area for improvement. Cooperative debate, where multiple agents generate, critique, and refine each other's responses, has been shown to improve factual accuracy and reasoning quality. In real-world negotiation, decisions are rarely made by a single individual; corporate acquisitions, policy agreements, and legal settlements typically involve teams that deliberate internally before engaging with the opposing side. This observation motivates the question of whether team-based negotiation — where each side is represented by a group of LLM agents that debate before responding — can lead to stronger negotiation outcomes.

This thesis investigates whether inference-time reasoning techniques and multi-agent collaboration can improve the negotiation capabilities of large language models. The work is organized around three main contributions:

\begin{itemize}
    \item How do open-source models compare to proprietary models in competitive negotiation scenarios?
    \item As inference-time reasoning techniques have shown promise individually, to what extent do they affect outcomes in multi-agent negotiation contexts?
    \item Can team-based negotiation, where each side is represented by a group of LLM agents that deliberate before responding, lead to stronger negotiation outcomes? 
\end{itemize}

The remainder of this document is organized as follows. Chapter 2 provides a literature review covering the evolution of language models, inference-time reasoning techniques, multi-agent systems, and the different interaction scenarios relevant to this work. Chapter 3 presents the detailed work plan for the upcoming months.